%! Polynomial Operations — Unit 3, Lesson 3.1 — Warm-Up
% The warm-up is set at 12pt (user direction, 2026-08-31) — larger type than the rest of the
% packet, deliberately — and it must still fit ONE page, blank and keyed. Three prerequisite
% items (Unit 3 has no 3.0): exponent rules, distributing, combining like terms. Item 3 ends on
% the question Round 1 answers: what a minus sign in front of parentheses does to EVERY term.
\documentclass[12pt]{article}
\usepackage{algebra2-article}
\usepackage{algebra2-boxes}

\begin{document}

\pageheader{Unit 3, Lesson 3.1}{Warm-Up}
% NAMESTRIP: no \namedateperiod here — the student writes their name once, on the cover.

\vspace{0.10in}

\begin{notesbox}{Getting Ready: Skills We'll Use Today}
Three things you already know. Today you add, subtract, and multiply whole polynomials, and
these are the moves underneath --- work quickly, then check with a neighbour.

\vspace{8pt}
\textbf{1. Exponent rules.} When you multiply, coefficients multiply and exponents on the same
variable add. Simplify each product.
\begin{center}
(a)\; $x^3\cdot x^4=$ \blank{1.6cm} \qquad
(b)\; $(-2x^2)(5x^3)=$ \blank{1.9cm} \qquad
(c)\; $(3x)(4xy)=$ \blank{1.9cm}
\end{center}

\vspace{8pt}
\textbf{2. Distribute.} Multiply the outside term by \emph{each} term inside.
\begin{center}
(a)\; $4x(x+2)=$ \blank{2.8cm} \qquad\qquad
(b)\; $-3(2x-5)=$ \blank{2.8cm}
\end{center}

\vspace{8pt}
\textbf{3. Combine like terms.} Like terms have the same variable raised to the same power.
\begin{itemize}[leftmargin=1.6em, itemsep=8pt]
  \item Simplify: \; $7x^2-3x+2-x^2+5x=$ \blank{3.4cm}
  \item A minus sign in front of parentheses means $-1$ times everything inside: \\[2pt]
        $-(x^2-4x+6)=$ \blank{3.4cm}
  \item How many of the three terms inside changed sign? \blank{3.2cm}
        (only the first / all three)
\end{itemize}
\textit{Hold that thought --- Round 1 subtracts a whole polynomial.}
\end{notesbox}

\end{document}
